\begin{center}
{\sffamily\Large\bfseries\color{purple}Earn the Next Turn}\par
\vspace{3mm}
{\large From Reusable Work to Accountable Succession}\par
\vspace{2mm}
{\small Vincent Boucher\quad $\cdot$\quad MONTREAL.AI}
\end{center}

\begin{insight}[The result, before the ambition]
\textbf{Constructive succession executes; the registered net-benefit gate remains on HOLD.}
The system retains a checked response instrument, uses it on new loads, and invokes it to construct a distinct
changed-network instrument. Eighteen such successors deliver 156 fresh, qualified parent-using outputs.
Against native refactorization with conventional caching, however, the whole-participant CPU ratio is
\textbf{0.9933}, with a paired-stream 95\% interval of \textbf{[0.9669, 1.0265]}. This does not support the
registered 5\% benefit. All seven treatments qualify all 492 offered useful outcomes.
\end{insight}

\subsection*{Abstract}

Can completed work create a method that improves later work, helps construct a different capability, and
pays for its own continuation? We introduce a reproducible causal succession experiment through the existing
AGI ALPHA workbench. A native sparse solve supplies a reusable factorization as a byproduct of useful service.
Fresh loads invoke that parent. Positive edge changes selected from the actual accepted predecessor then
produce a low-rank response instrument through parent solves. Every output receives a separately implemented
edge-balance and energy check; descendants inherit neither current proof status nor operational authority.

A source and protocol freeze precedes one confirmation campaign of 24 independent network streams, seven
treatments and four workload/maintenance regimes. Controls retain ordinary caching, warm starts, competent
native construction, an independently built equivalent parent, and equally enabled conventional reuse.
All 3,444 offered useful outputs qualify. The flywheel constructs eighteen successors and qualifies 156 fresh
outputs that execute their parent, improving on the prior edition's unactivated productive lineage. But the
primary whole-participant CPU test fails: the point estimate is only 0.67\% below native refactorization and its
interval includes a loss. Conventional reuse satisfies the registered noninferiority guard, without evidence
of unique agentic advantage. Actual released service quota funds 1,034.057\,ms of successor dispatch; the quota
originates in external endowment and does not prove self-financing growth.

The contribution is an executed, falsifiable boundary between useful inheritance, construction dependence,
whole-system value and resource feedback. An offline browser companion, exact-source native implementation,
protected confirmation record, unsuccessful development attempts, raw data and independent reanalysis are
provided. Full research investment recovery, commercial growth and acceleration remain unestablished.

\vspace{2mm}
{\small\textbf{Keywords:} executable inheritance; constructive succession; sparse numerical instruments;
full-cost evaluation; resource custody; bounded autonomy; reproducibility.}

\subsection*{What the evidence establishes}

\begin{tabularx}{\linewidth}{@{}L{44mm}Y@{}}\toprule
Question & Finding \\\midrule
Does the recurring loop run? & Yes. One launch executes 168 finite trials without manual phase triggers. \\
Does accepted work leave useful means? & Yes. Fresh solves use retained factors; actual B08 selects the successor's changed operator. \\
Does the parent construct a distinct child? & Yes. Parent calls produce the response basis; fresh changed-network outputs qualify. \\
Does net benefit survive the strongest comparison? & Not established. The registered 5\% whole-participant CPU gate fails. \\
Does released capacity fund construction? & Yes, within a conserved externally supplied quota; incremental earnings are not proved. \\
Does the system accelerate or recover all investment? & Not established. No recursive discovery-rate experiment or complete monetary/human-cost evidence. \\\bottomrule
\end{tabularx}
